\documentclass{ximera}

\begin{document}

    \title{Homework 6}
    
    \begin{abstract} Solve the following problems.
    \\
    Due Date: Sunday, November 15th, 11:59 pm
    \end{abstract}
    \maketitle

\section{Projections}
\begin{exercise}
Let
\[
A = \begin{pmatrix}
1 & 0\\
0 & 0
\end{pmatrix}.
\]
    
\begin{enumerate}
    \item Compute $A^2$ and verify that $A^2=A$.
    \item Explain why $A$ is a projection matrix.
    \item For the vector
    \[
    x = \begin{pmatrix} 3\\ 4 \end{pmatrix},
    \]
    compute $Ax$. What is the geometric interpretation of this projection?
\end{enumerate}

\end{exercise}


\begin{exercise}
Let
\[
A = \frac{1}{2}
\begin{pmatrix}
1 & 1\\
1 & 1
\end{pmatrix}.
\]

\begin{enumerate}
    \item Verify that $A^2=A$.
    \item Compute $Ax$ for
    \[
    x = \begin{pmatrix} 2\\ 4 \end{pmatrix}.
    \]
    \item Describe geometrically the subspace onto which $A$ projects.
\end{enumerate}

\end{exercise}

\begin{exercise} (\textbf{no grade} geometry review)
    Recall the the formula for a projection is $P_w(v) = \frac{v\cdot w}{w \cdot w} w$.  Assume $v,w \in \mathbb{R}^2$.  Use geometry to show that $v\cdot w = \|v\|_2\|w\|_2 \cos(\theta)$, where $\theta$ is the angle between $v$ and $w$.  (This geometry formula holds also for vectors in $\mathbb{R}^n$.)
\end{exercise}

\section{LU Decomposition}
% \begin{exercise}
% Find an LU decomposition of the matrix
% \[
% A = \begin{pmatrix}
% 2 & 1\\
% 4 & 3
% \end{pmatrix},
% \]
% where $L$ is lower triangular with ones on the diagonal and $U$ is upper triangular.

% Verify your answer by computing $LU$.

% \end{exercise}


\begin{exercise}
Use LU decomposition to solve the linear system
\[
\begin{aligned}
2x+y &= 5,\\
4x+3y &= 11.
\end{aligned}
\]

In particular:

\begin{enumerate}
    \item Write the system in matrix form $Ax=b$.
    \item Find an LU decomposition of $A$.
    \item Solve $Lc=b$ using forward substitution.
    \item Solve $Ux=c$ using back substitution.
\end{enumerate}

\end{exercise}

\begin{exercise}
Consider the matrix
\[
A = \begin{pmatrix}
10^{-4} & 1\\
1 & 1
\end{pmatrix}.
\]

\begin{enumerate}
    \item Perform Gaussian elimination without pivoting to obtain an upper triangular matrix $U$. What is the multiplier used in the elimination step?

    \item Now perform Gaussian elimination with partial pivoting. What row exchange is performed, and what is the resulting upper triangular matrix?

    \item Explain why partial pivoting is preferable in this example. In particular, compare the sizes of the multipliers obtained in the two approaches.
\end{enumerate}

\textit{A large multiplier does not automatically mean that the final answer will be inaccurate, but it can lead to significant error accumulation in floating-point arithmetic.}

\end{exercise}


\end{document}
